\section{Limitations}
\label{sec:limitations}

The evidence has five explicit boundaries. First, the main scenarios are generated on a straight road with constant-velocity obstacle models. This may overstate the regularity of the predicted occupancy and does not test curvature, lane topology, weather, or incident effects; the next study should replay public and field-collected trajectories across road geometries. Second, the prediction-noise model uses bounded position and velocity errors. Its collision result therefore applies to the stated error envelope, not to an arbitrary predictor; calibrated multi-modal prediction distributions are needed next. Third, B0--B2 are internal matched PVD variants and B3 is a finite coarse-grid reference. The paper makes no public SOTA or fair cross-dataset claim; a matched implementation of a recent external planner is the next comparison. Fourth, Apollo evidence covers interface compatibility and representative Dreamview behavior, not native Apollo performance, hardware-in-the-loop execution, physical-road testing, or an end-to-end safety proof. Native closed-loop instrumentation and controlled-road tests are required before deployment claims. Fifth, the maximum-gap theorem is exact for a fixed-section interval abstraction and the retained candidate set. It does not certify a continuous global optimum of the full nonconvex trajectory problem; a future formulation could study completeness under changing obstacle order.

\section{Conclusion}
\label{sec:conclusion}

This paper presented MIKU, an interaction-aware constraint compiler for path--velocity planning with dynamic multi-obstacle traffic. Arrival-time-dependent projection, continuous group-level partitioning, uncertainty-expanded occupancy, and temporal homotopy convert lost interaction information into SL and ST bounds while preserving the downstream path and speed QPs. The fixed-section theorem reduces side-assignment search to an $O(k\log k)$ scan, which matched exhaustive enumeration on 4,000 deterministic instances. In 3,500 paired scenarios, collision-free goal reaching increased from 62.7\% to 77.5\% and collision decreased from 2.2\% to 1.1\% relative to the time-blind baseline, with the largest effects in narrow and delayed-crossing configurations. Typical latency remained near 10 ms, but the P99 tail was higher and rolling replanning increased episode-runtime P95, so the result supports bounded simulation and interface claims rather than a hard real-time or road-deployment guarantee. The next study should validate the same constraint transfer with calibrated public or field trajectories and native runtime instrumentation.
